\documentclass[10pt,a4paper]{amsart}
\usepackage{amsmath,amssymb,amsthm,graphicx}
\usepackage{verbatim}
\newenvironment{reference}{\comment}{\endcomment}
%\newenvironment{reference}{}{}
\newenvironment{slogan}{\comment}{\endcomment}
\newenvironment{history}{\comment}{\endcomment}

% For commutative diagrams we use Xy-pic
\usepackage[all]{xy}

% We use 2cell for 2-commutative diagrams.
\xyoption{2cell}
% Theorem environments.
%
\theoremstyle{plain}
\newtheorem{theorem}[subsection]{Theorem}
\newtheorem{proposition}[subsection]{Proposition}
\newtheorem{lemma}[subsection]{Lemma}

\theoremstyle{definition}
\newtheorem{definition}[subsection]{Definition}
\newtheorem{example}[subsection]{Example}
\newtheorem{exercise}[subsection]{Exercise}
\newtheorem{situation}[subsection]{Situation}

\theoremstyle{remark}
\newtheorem{remark}[subsection]{Remark}
\newtheorem{remarks}[subsection]{Remarks}

\numberwithin{equation}{subsection}

% Macros
%
\def\lim{\mathop{\mathrm{lim}}\nolimits}
\def\colim{\mathop{\mathrm{colim}}\nolimits}
\def\Spec{\mathop{\mathrm{Spec}}}
\def\Hom{\mathop{\mathrm{Hom}}\nolimits}
\def\Ext{\mathop{\mathrm{Ext}}\nolimits}
\def\SheafHom{\mathop{\mathcal{H}\!\mathit{om}}\nolimits}
\def\SheafExt{\mathop{\mathcal{E}\!\mathit{xt}}\nolimits}
\def\Sch{\mathit{Sch}}
\def\Mor{\mathop{\mathrm{Mor}}\nolimits}
\def\Ob{\mathop{\mathrm{Ob}}\nolimits}
\def\Sh{\mathop{\mathit{Sh}}\nolimits}
\def\NL{\mathop{N\!L}\nolimits}
\def\CH{\mathop{\mathrm{CH}}\nolimits}
\def\proetale{{pro\text{-}\acute{e}tale}}
\def\etale{{\acute{e}tale}}
\def\QCoh{\mathit{QCoh}}
\def\Ker{\mathop{\mathrm{Ker}}}
\def\Im{\mathop{\mathrm{Im}}}
\def\Coker{\mathop{\mathrm{Coker}}}
\def\Coim{\mathop{\mathrm{Coim}}}

% Boxtimes
%
\DeclareMathSymbol{\boxtimes}{\mathbin}{AMSa}{"02}

%
% Macros for moduli stacks/spaces
%
\def\QCohstack{\mathcal{QC}\!\mathit{oh}}
\def\Cohstack{\mathcal{C}\!\mathit{oh}}
\def\Spacesstack{\mathcal{S}\!\mathit{paces}}
\def\Quotfunctor{\mathrm{Quot}}
\def\Hilbfunctor{\mathrm{Hilb}}
\def\Curvesstack{\mathcal{C}\!\mathit{urves}}
\def\Polarizedstack{\mathcal{P}\!\mathit{olarized}}
\def\Complexesstack{\mathcal{C}\!\mathit{omplexes}}
% \Pic is the operator that assigns to X its picard group, usage \Pic(X)
% \Picardstack_{X/B} denotes the Picard stack of X over B
% \Picardfunctor_{X/B} denotes the Picard functor of X over B
\def\Pic{\mathop{\mathrm{Pic}}\nolimits}
\def\Picardstack{\mathcal{P}\!\mathit{ic}}
\def\Picardfunctor{\mathrm{Pic}}
\def\Deformationcategory{\mathcal{D}\!\mathit{ef}}
\usepackage{cleveref}
\crefname{originallemma}{original lemma}{original lemmas}
\Crefname{originallemma}{Original lemma}{Original lemmas}
\makeatletter
\newlabel{lemma-shriek-closed-immersion}{{0AA2}{source}{}{sourcepointer-1}{}}
\newlabel{lemma-shriek-closed-immersion@cref}{{[originallemma][0][]0AA2}{[1][0][]source}{}{}{}}
\makeatother
\crefname{originallemma}{original lemma}{original lemmas}
\Crefname{originallemma}{Original lemma}{Original lemmas}
\makeatletter
\newlabel{morphisms-lemma-closed-point-fibre-locally-finite-type}{{01TF}{source}{}{sourcepointer-2}{}}
\newlabel{morphisms-lemma-closed-point-fibre-locally-finite-type@cref}{{[originallemma][0][]01TF}{[1][0][]source}{}{}{}}
\makeatother
\crefname{originalsection}{original section}{original sections}
\Crefname{originalsection}{Original section}{Original sections}
\makeatletter
\newlabel{example-iso-on-RSheafHom-noetherian}{{0GEW}{source}{}{sourcepointer-3}{}}
\newlabel{example-iso-on-RSheafHom-noetherian@cref}{{[originalsection][0][]0GEW}{[1][0][]source}{}{}{}}
\makeatother
\title{The normalised dualizing complex is normalised at closed points lying over closed points}
\author{The Stacks Project Authors}
\date{stacks-project:duality.tex:local-snapshot-20261012; GNU FDL 1.2 or later}
\begin{document}
\maketitle
\noindent\textit{Editorial note.} All delivered mathematical text is the author's own, copied from the pinned Stacks Project chapter file duality.tex; no statement or proof step was rewritten, reordered or removed.\par\medskip
\noindent\textit{Editorial note.} Preamble provenance: the pinned chapter file opens with an input of the shared preamble and carries no preamble text of its own, so the project's single Stacks-wide preamble staged at sources/stacks-duality/preamble.tex is excerpted verbatim under the recipe fixed by LANE\_BRIEF 22.1, as three spans, lines 8-13 (verbatim package and the reference/slogan/history environments), lines 14-18 (Xy-pic setup) and lines 154-217 (the native theorem declarations and mathematical macros). Nothing was added, removed or redefined; the wrapper class and the amsmath/amssymb/amsthm/graphicx base packages are supplied by the assembler.\par\medskip
\noindent\textit{Editorial note.} The counted statement and proof contain no cross reference, no citation and no included file, so the delivered body is self contained.\par\medskip
\noindent\textit{Editorial note.} Canonical location of the counted proposition: Stacks tag 0AWL (sources/stacks-duality/pinned-original-source.tex line 5805).\par\medskip
\noindent\textit{Allowed standard prerequisites (editorial).}\begin{itemize}
\item The upper shriek of a closed immersion and its compatibility with composition. Same chapter, Lemma labelled lemma-shriek-closed-immersion, cited at proof lines 5826-5830. Identifies the localised Hom with i\_x\textasciicircum{}! of the dualizing complex.
\item Finiteness of the residue field extension at a closed point of a finite type morphism. Morphisms, Lemma labelled morphisms-lemma-closed-point-fibre-locally-finite-type, cited at proof lines 5836-5840. Shows pi is finite, so that the composition of shrieks applies.
\item The comparison between pushforward and internal Hom for a Noetherian morphism. Same chapter, Example labelled example-iso-on-RSheafHom-noetherian, cited at proof lines 5850-5852. Gives the first equality computing Rpi\_* of the shriek.
\item Normalised dualizing complexes and their description through the residue field. Same chapter, the definition of a normalised dualizing complex, invoked in the statement and used at proof lines 5842-5846. Turns normalisation into i\_s\textasciicircum{}!omega\_S = kappa(s)[0].
\end{itemize}\medskip
\ifdefined\hypertarget\hypertarget{sourcepointer-1}{}\fi
\noindent\textit{Source reference.} Original lemma 0AA2: sources/stacks-duality/pinned-original-source.tex, source lines 4624--4624. Same-chapter pointer into the author environment lemma at source line 4624; invoked at proof line 5825 of the counted proof. It is retained only as a navigation pointer; every load-bearing step is inside the delivered span.\par\medskip
\ifdefined\hypertarget\hypertarget{sourcepointer-2}{}\fi
\noindent\textit{Source reference.} Original lemma 01TF: sources/stacks-chow/tags.txt, source lines 2372--2372. Cross-chapter standard prerequisite located in the project tag list staged in the local snapshot; invoked at proof line 5834 of the counted proof. It is retained only as a navigation pointer; every load-bearing step is inside the delivered span.\par\medskip
\ifdefined\hypertarget\hypertarget{sourcepointer-3}{}\fi
\noindent\textit{Source reference.} Original section 0GEW: sources/stacks-duality/pinned-original-source.tex, source lines 865--865. Same-chapter pointer into the author environment example at source line 865; invoked at proof line 5849 of the counted proof. It is retained only as a navigation pointer; every load-bearing step is inside the delivered span.\par\medskip
\begin{lemma}
\label{lemma-good-dualizing-normalized}
Let $S$ be a Noetherian scheme and let $\omega_S^\bullet$ be a
dualizing complex. Let $X$ be a scheme of finite type over $S$ and let
$\omega_X^\bullet$ be the dualizing complex normalized relative
to $\omega_S^\bullet$. If $x \in X$ is a closed point lying over
a closed point $s$ of $S$, then $\omega_{X, x}^\bullet$
is a normalized dualizing complex over $\mathcal{O}_{X, x}$
provided that $\omega_{S, s}^\bullet$ is a normalized dualizing
complex over $\mathcal{O}_{S, s}$.
\end{lemma}

\begin{proof}
We may replace $X$ by an affine neighbourhood of $x$, hence we may
and do assume that $f : X \to S$ is separated.
Then $\omega_X^\bullet = f^!\omega_S^\bullet$. We have to show that
$R\Hom_{\mathcal{O}_{X, x}}(\kappa(x), \omega_{X, x}^\bullet)$
is sitting in degree $0$. Let $i_x : x \to X$ denote the inclusion
morphism which is a closed immersion as $x$ is a closed point.
Hence $R\Hom_{\mathcal{O}_{X, x}}(\kappa(x), \omega_{X, x}^\bullet)$
represents $i_x^!\omega_X^\bullet$ by
Lemma \ref{lemma-shriek-closed-immersion}.
Consider the commutative diagram
$$
\xymatrix{
x \ar[r]_{i_x} \ar[d]_\pi & X \ar[d]^f \\
s \ar[r]^{i_s} & S
}
$$
By Morphisms, Lemma
\ref{morphisms-lemma-closed-point-fibre-locally-finite-type}
the extension $\kappa(x)/\kappa(s)$ is finite and hence
$\pi$ is a finite morphism. We conclude that
$$
i_x^!\omega_X^\bullet = i_x^! f^! \omega_S^\bullet =
\pi^! i_s^! \omega_S^\bullet
$$
Thus if $\omega_{S, s}^\bullet$ is a normalized dualizing complex
over $\mathcal{O}_{S, s}$, then $i_s^!\omega_S^\bullet = \kappa(s)[0]$
by the same reasoning as above. We have
$$
R\pi_*(\pi^!(\kappa(s)[0])) =
R\SheafHom_{\mathcal{O}_s}(R\pi_*(\kappa(x)[0]), \kappa(s)[0]) =
\widetilde{\Hom_{\kappa(s)}(\kappa(x), \kappa(s))}
$$
The first equality by Example \ref{example-iso-on-RSheafHom-noetherian}
applied with $L = \kappa(x)[0]$. The second equality holds because
$\pi_*$ is exact.
Thus $\pi^!(\kappa(s)[0])$ is supported in degree $0$ and we win.
\end{proof}

\end{document}
